\subsection{The marked cubic orbifold over the nine-point translation cover}

\begin{result}{A.4.44}{Proved}{The order-twenty-seven source and its simultaneous cone and cusp detector}
There is an actual orbifold extension
\[
 1\longrightarrow\Pi_9\longrightarrow
 E=\langle u,v\mid u^3=v^3=1\rangle
 \longrightarrow P=H\rtimes C_3\longrightarrow1,
 \qquad H\cong C_3^2,
                                                               \tag{COB.1}
\]
for the torus punctured at its nine three-division points.
Here $P$ is the nonabelian group of order $27$ and
exponent three. The intermediate once-punctured torus
group and its peripheral word have the literal marking
\[
 \begin{gathered}
 \Sigma=\ker\bigl(E\to C_3,\ u,v\mapsto1\bigr)
       =F(x,y),\qquad
 x=vu^{-1},\quad y=uvu^{-2},\\
 \Pi_9=\Sigma'\Sigma^3,\qquad
 c=(uv)^{-1},\qquad c^3=[x,y],\\
 uxu^{-1}=y,\qquad uyu^{-1}=(xy)^{-1},\qquad
 R_0=\begin{pmatrix}0&-1\\1&-1\end{pmatrix}.
 \end{gathered}
                                                               \tag{COB.2}
\]
The marked linear CM element is the image of the
infinite-order cusp word $c$. The two finite factors
$\langle u\rangle,\langle v\rangle$ are the other
two affine CM classes; their fixed points are unmarked.

Let $q:E\twoheadrightarrow D$ be any specified finite
quotient through which the map to $P$ in (COB.1)
factors. Put $K=\ker q$, $n=|D|$, and
\[
 D^+=E/(K'K^3),\qquad
 |D^+|=n3^{1+n/3}.
                                                               \tag{COB.3}
\]
For each $t\in\{u,v,c\}$, the exact centralizer
projection is
\[
 \operatorname{im}\bigl(C_{D^+}(t)\longrightarrow D\bigr)
                  =\langle q(t)\rangle.
                                                               \tag{COB.4}
\]
If $m$ is the order of $q(c)$, the order of $c$ in
$D^+$ is exactly $3m$. For the particular quotient
$D=P$, (COB.3) has order $3^{13}$; all $27$ lower
residues have been computed. Each of the three
centralizer fibres in (COB.4) has kernel dimension
four over $\mathbb F_3$. Exactly twelve nonidentity
residues admit an order-three lift, the two signed
cone types and their conjugates. The marked cusp
residues do not admit such a lift.

These assertions are about the displayed orbifold
group and its actual finite quotients. They do not
identify the normalizer of $H$ in a finite ambient
nine-point mapping class group with the image of its
ordinary normalizer.
\end{result}

\begin{proof}
\par\medskip\noindent\textbf{The free-group marking.}\quad
Use the three coset representatives $1,u,u^2$ for the
kernel of $u,v\mapsto1$ in $C_3$. Its three nontrivial
Schreier entries before eliminating the relator are
\[
 x=vu^{-1},\quad y=uvu^{-2},\quad z=u^2v,
                 \qquad xyz=1.
                                                               \tag{COB.5}
\]
The relator is the rewrite of $v^3$; $u^3$ eliminates
the remaining tree entry. Thus $\Sigma=F(x,y)$,
and conjugation by $u$ cycles $x,y,z$. This proves
the two automorphism formulas in (COB.2), with their
actual abelianization matrix $R_0$. Direct substitution
also gives
\[
                  (uv)^3=yxy^{-1}x^{-1}=[y,x].
                                                               \tag{COB.6}
\]
It proves the precise peripheral identity in (COB.2),
including its orientation.

For a literal finite model of $P$ use
\[
 \begin{gathered}
 (i,j,k)(I,J,K)=(i+I,j+J,k+K+iJ)\quad\text{in }\mathbb F_3^3,\\
 u=(1,0,0),\qquad v=(0,1,0).
 \end{gathered}
                                                               \tag{COB.7}
\]
The elements $u,v$ generate all $27$ triples, and every
triple has cube one. Their commutator $(0,0,1)$ is
nonidentity. The intermediate subgroup is
$H=\{(i,j,k):i+j=0\}$, and
\[
 x=(2,1,0),\qquad y=(2,1,1),\qquad c=(2,2,0).
                                                               \tag{COB.8}
\]
The nine products $x^a y^b$ are distinct and exhaust
$H$. Therefore the kernel on $\Sigma$ is exactly
$\Sigma'\Sigma^3$, proving every group map in
(COB.1)--(COB.2). In particular its free rank is ten,
by the index-nine Schreier formula.

\par\medskip\noindent\textbf{The actual fixed-point marking.}\quad
Take the oriented hexagonal torus with lattice
coordinates in $\mathbb R^2/\mathbb Z^2$ and rotation
$R_0$. Its translations by $\frac13\mathbb Z^2/
\mathbb Z^2$ form $H$. The following two actual
affine transformations give the map (COB.7):
\[
 u(z)=R_0z+\tfrac13(1,0),\qquad
 v(z)=R_0z+\tfrac13(1,1),\qquad c(z)=R_0z.
                                                               \tag{COB.9}
\]
Indeed $uv=R_0^2$, and the two translations from
$x,y$ are respectively $(0,1)/3$ and $(-1,-1)/3$.
Their integral lifts $b_1=(0,1)$ and $b_2=(-1,-1)$
have determinant one. The upstairs lattice is
$\Lambda=\mathbb Zb_1+\mathbb Zb_2$; the quotient
torus has lattice $\frac13\Lambda$ and primitive
basis $b_1/3,b_2/3$. The covering sends each upstairs
$b_i$ to $3(b_i/3)$, so its period map is the actual
matrix $3I$. No division by three after finite
reduction is used.

This affine marking can also be checked before taking
any finite quotient. Impose $c^3=1$ in $E$.
By (COB.2) this imposes exactly $[x,y]=1$ on
$\Sigma$, giving the group $\mathbb Z^2\rtimes C_3$
with the displayed rotation $R_0$. It acts faithfully
on the oriented plane by translations $b_1,b_2$
and the two rotations with translation parts $(1,0)$
and $(1,1)$. Dividing this translation lattice by
$3\Lambda$ and rescaling the plane by three gives
exactly (COB.9). Thus the affine generators come
from the same orbifold presentation and marking;
they are not an identification chosen only from
the abstract order of $P$.

Solving $(I-R_0)z=t/3$ uses
\[
 (I-R_0)^{-1}=\tfrac13
                 \begin{pmatrix}2&-1\\1&1\end{pmatrix}.
                                                               \tag{COB.10}
\]
The fixed triples, recorded by their numerators in
$\frac19\mathbb Z^2/\mathbb Z^2$, are exactly
\[
\begin{array}{c|c|c}
 &\text{fixed triple}&\text{points in }H\\ \hline
 u&(2,1),(5,7),(8,4)&0\\
 v&(1,2),(4,8),(7,5)&0\\
 c&(0,0),(3,6),(6,3)&3.
\end{array}
                                                               \tag{COB.11}
\]
Each full $P$-orbit of these triples has nine points:
it is its translation orbit under $H$. The three
orbits are disjoint. Every nontranslation element of
$P$ has three fixed points, because $I-R_0$ and
$I-R_0^2$ both have determinant three. There are
eighteen such elements, giving $54$ nonidentity
stabilizer incidences. The displayed $27$ points,
each with two nonidentity stabilizer elements, already
account for all of them.

The quotient of the compact torus by $P$ is a sphere.
Indeed its ramification formula is
$0=27(2-2g)-27(3-1)$, giving $g=0$.
Its three branch orbits have stabilizer order three.
Removing precisely the orbit $H$ removes the third
branch point and leaves two order-three cone points.
The resulting oriented orbifold has fundamental group
$C_3*C_3$. Positive loops about its two cone points
are $u,v$, and the positive puncture loop is
$c=(uv)^{-1}$. This identifies the topological extension
with (COB.1), including the named puncture and the two
unmarked fixed-point orbits. The intermediate quotient
by $H$ is the once-punctured torus with group $\Sigma$.
In particular the linear CM class must not be
substituted for one of the two cone generators:
its actual lift in this extension is $c$, with the
nontrivial third power in (COB.2).

\par\medskip\noindent\textbf{One relative graph layer.}\quad
Since $q$ retains $P$, it is injective on both finite
factors. The Bass--Serre quotient graph for $K$ has
vertices $D/\langle q(u)\rangle$ and
$D/\langle q(v)\rangle$, and oriented edges $D$.
The edge $g$ goes from $g\langle q(u)\rangle$ to
$g\langle q(v)\rangle$. It is connected, has
$2n/3$ vertices and $n$ edges, and hence has free
fundamental group of rank $1+n/3$. This proves
(COB.3) and identifies its kernel with the actual
graph homology over $\mathbb F_3$.

For the two cone generators, (COB.4) follows directly
from the invariant-chain proof of A.4.7. In this graph
an invariant edge chain has boundary coefficient zero
at every fixed vertex. A commuting lift moving the
chosen fixed vertex would give such a chain with
nonzero boundary there. Thus its lower value fixes
that vertex and lies in the corresponding cyclic
factor. Those cyclic values have the stated literal
commuting lifts.

For the cusp, the graph-chain calculation is equally
explicit. Write $e_g$ for the edge $g$. From the base
vertex of type $u$, choose paths with crossed chain
\[
 D(u)=0,\qquad D(v)=e_1-e_{q(v)},\qquad
 D(c)=e_1-e_{q(v)^{-1}},\qquad
 D(ab)=D(a)+q(a)D(b).
                                                               \tag{COB.12}
\]
The two relators $u^3,v^3$ have zero chain, so these
are literal chains of the specified orbifold words.
With $m=\operatorname{ord}(q(c))$, the element $c^m$
lies in $K$ and has homology chain
\[
        w=\sum_{j=0}^{m-1}e_{q(c)^j}
               -\sum_{j=0}^{m-1}e_{q(c)^jq(v)^{-1}}.
                                                               \tag{COB.13}
\]
The two supports are distinct right cosets of
$\langle q(c)\rangle$: otherwise $q(v)$ belongs to
that cyclic group, contradicting its retained image
in the nonabelian group (COB.7). Thus $w\ne0$.
Since $1\ne-1$ in $\mathbb F_3$, every deck element
fixing $w$ preserves its positive support. The exact
stabilizer of that support under left translation is
$\langle q(c)\rangle$.

A lift centralizing $c$ fixes $c^m$, so its lower
deck value fixes $w$ and belongs to this cyclic
subgroup. Every cyclic value is represented by a
power of $c$, proving (COB.4) for the third element.
The nonzero class $w$ has order three, so $c$ has
order exactly $3m$ in $D^+$.

\par\medskip\noindent\textbf{The complete first-layer certificate.}\quad
For $D=P$, retain all $27$ triples in lexicographic
order. The program
\path{work/tower_code/cubic_orbit_orbifold.py}
constructs the $18\times27$ incidence matrix, of
rank $17$, and a $27\times10$ basis $B$ of its
kernel. The crossed-chain map
$E\to\mathbb F_3^{27}\rtimes P$ from (COB.12)
has kernel exactly $\Pi_9'\Pi_9^3$: on $\Pi_9$
it is the inclusion of the graph's first homology
in its edge chains. All $54$ Schreier edges are
recorded and span the ten-dimensional kernel.
Thus this is a faithful finite representation of
the actual group of order $3^{13}$, rather than
an ambient space of formal candidate vectors.

For each $h\in P$ let $D_h$ be the chain of its
recorded source-word section. A lift of $h$ has
exactly the chains $D_h+Bz$. Commuting with
$t\in\{u,v,c\}$ first requires $h q(t)=q(t)h$,
and then the exact linear equation
\[
 (q(t)-1)Bz=(h-1)D(t)-(q(t)-1)D_h.
                                                               \tag{COB.14}
\]
The program solves all $27$ residues for each of
the three elements. Each matrix on the left has
rank six. Precisely the three residues in (COB.4)
survive, each with four free coefficients.
The cube-one equation is likewise the literal
norm equation
\[
 (1+h+h^2)Bz=-(1+h+h^2)D_h.
                                                               \tag{COB.15}
\]
All $27$ tests are retained. Its twelve nonidentity
survivors are exactly the conjugates of
$u,u^{-1},v,v^{-1}$; each has a literal cone lift.
The exclusion of every other nonidentity residue
also follows from the order-three recognition
argument of A.4.7. Finally the program verifies
the nonzero six-edge chain of $c^3$, its exact
three-element stabilizer, all affine maps (COB.9),
and all fixed points in (COB.11).
Its output, JSON record, and matrix archive retain
every equation and source section.

The quotient with $D=P$ is solvable and so does
not remove the A.4.42 discrepancy by adjoining it
as another coordinate. Formula (COB.3), however,
is a specified relative graph layer for every
retained finite orbifold input $D$. Inserting such
a layer into the ambient mapping-class recognition
still requires the actual common surface input,
its geometric action, and the comparison with the
eight-point filling kernel. Those maps have not
been inferred from this cyclic computation.
\end{proof}
